% !TeX root = tcolorbox.tex
% include file of tcolorbox.tex (manual of the LaTeX package tcolorbox)
\clearpage
\section{简介}%
\tcbset{external/prefix=external/intro_}%
本宏包源于 2006 年前后拙作
\flqq\citetitle{sturm:latex}\frqq~\cite{sturm:latex}
第一版的编写过程。
为了在书中同时展示 \LaTeX\ 源代码和编译效果，
我希望用醒目的彩色文本框来并排呈现二者。
后来我发现这种文本框也非常适合用来突出定义和定理，
于是也将其用于我的数学讲义
\cite{sturm:mathe1,sturm:mathe2,sturm:mathe3}。
通过本宏包，欢迎你在类似的项目中使用这些文本框。

2.00 版本的重大更新是支持了可跨页断开的文本框。
这一特性在不影响核心宏包的前提下扩展了应用场景——
如果你不需要文本框自动跨页，完全可以忽略它。
2.20 版本加入了呼声很高的代码列表"并排"模式。
3.00 版本引入了带框标题，并大幅增强了
覆盖层、底层修饰、收尾装饰以及自定义代码扩展的定制选项。

\begin{tcolorbox}[enhanced,boxrule=0mm,boxsep=0mm,frame hidden,interior hidden,
  left=0mm,right=0mm,top=0mm,bottom=0mm,watermark opacity=0.25,watermark zoom=1.2,before=\par\smallskip,
  clip watermark=false,
  watermark tikz={%
    \path[fill=yellow,draw=yellow!75!red] (0,0) circle (1cm);
    \fill[red] (45:5mm) circle (1mm);
    \fill[red] (135:5mm) circle (1mm);
    \draw[line width=1mm,red] (215:5mm) arc (215:325:5mm);}]
自 2011 年首次公开发布以来，我收到了来自世界各地的大量反馈。
在此感谢所有来信告知问题、提出新功能或改进建议的朋友，
正是你们的支持让本宏包不断完善。
\end{tcolorbox}


\subsection{安装}
通常情况下，|tcolorbox| 会作为主流 \LaTeX\ 发行套件的一部分被安装，
用户无需额外操作。

如果你打算\emph{手动}本地安装，
请参考 |tcolorbox| 宏包的 |README| 文件中的说明。
简言之：你不仅要安装 |tcolorbox.sty|，
还必须将所有 |*.code.tex| 文件一起安装到本地 |texmf| 树中。


\subsection{加载宏包}
核心宏包 |tcolorbox| 会自动加载以下宏包：
\refPkg{tikz} \cite{tantau:tikz_and_pgf}、
\refPkg{verbatim} \cite{schoepf:verbatim}、
\refPkg{etoolbox} \cite{lehmann:etoolbox}、
以及 \refPkg{environ} \cite{robertson:2014a}。
在导言区中按常规方式加载 |tcolorbox|：
\begin{dispListing}
\usepackage{tcolorbox}
\end{dispListing}
该宏包采用键值语法接收选项。
你也可以在导言区后续使用
\refCom{tcbuselibrary}（参见该处）来加载选项。
例如，排版代码列表的选项为：
\begin{dispListing}
\usepackage[listings]{tcolorbox}
\end{dispListing}


\clearpage
\subsection{程序库}\label{sec:bibliothek}
核心宏包 |tcolorbox| 可通过程序库进行扩展。
加载宏包时使用选项键，或在导言区中使用以下命令
并传入相同的键集合即可。

\begin{docCommand}{tcbuselibrary}{\marg{key list}}
  加载由 \meta{键列表} 指定的程序库。
\begin{dispListing}
\tcbuselibrary{listings,theorems}
\end{dispListing}
\end{docCommand}

以下键可在 |\tcbuselibrary| 或 |\usepackage| 中使用，
无需写出完整键路径 |/tcb/library/|。

\begin{docTcbKey}[library]{skins}{}{\mylib{skins}}
  加载 \refPkg[tikzfill]{tikzfill.image} 宏包，
  并为彩色文本框的外观提供额外的样式（皮肤）；
  参见 \zvref{sec:skins}。
\end{docTcbKey}

\begin{docTcbKey}[library]{vignette}{}{\mylib{vignette}}
  提供更多装饰性图案代码；参见 \zvref{sec:vignette}。
\end{docTcbKey}

\begin{docTcbKey}[library]{raster}{}{\mylib{raster}}
  提供额外的命令和选项，用于将多个文本框排成网格状布局；
  参见 \zvref{sec:raster}。
\end{docTcbKey}

\begin{docTcbKey}[library]{listings}{}{\mylib{listings}}
  加载 \refPkg{listings} 宏包 \cite{hoffmann:listings}，
  并提供用于排版代码列表的额外命令，详见 \zvref{sec:listings}。
\end{docTcbKey}

\begin{docTcbKey}[library]{listingsutf8}{}{\mylib{listingsutf8}}
  加载 \refPkg{listings} \cite{hoffmann:listings} 和
  \refPkg{listingsutf8} \cite{oberdiek:listingsutf8} 宏包以支持 UTF-8。
  这是 \mylib{listings} 程序库的一个变体，
  详见 \zvref{sec:listings}。
\end{docTcbKey}

\begin{docTcbKey}[library]{minted}{}{\mylib{minted}}
  加载 \refPkg{minted} 宏包 \cite{poore:minted}，
  使用 |Pygments| \cite{pygments:web} 工具排版代码列表，
  另见 \zvref{sec:listings}。
\end{docTcbKey}

\begin{docTcbKey}[library]{theorems}{}{\mylib{theorems}}
  提供用于排版定理的额外命令，详见 \zvref{sec:theorems}。
\end{docTcbKey}


\begin{docTcbKey}[library]{breakable}{}{\mylib{breakable}}
  支持文本框自动跨页断开；参见 \zvref{sec:breakable}。
\end{docTcbKey}

\begin{docTcbKey}[library]{magazine}{}{\mylib{magazine}}
  支持将断开的文本框片段暂存，以便后续按原序或交换顺序使用，
  参见 \zvref{sec:magazine}。
\end{docTcbKey}

\begin{docTcbKey}[library]{poster}{}{\mylib{poster}}
  支持创建海报；参见 \zvref{sec:poster}。
\end{docTcbKey}

\begin{docTcbKey}[library]{fitting}{}{\mylib{fitting}}
  支持根据文本框尺寸自动调整框内字体大小；
  参见 \zvref{sec:fitting}。
\end{docTcbKey}

\begin{docTcbKey}[library]{hooks}{}{\mylib{hooks}}
  将多个选项键扩展为"可挂钩"的键；参见 \zvref{sec:hooks}。
\end{docTcbKey}

\clearpage
\begin{docTcbKey}[library]{xparse}{}{\mylib{xparse}}
  加载 \refPkg{xparse} 宏包，是为兼容性保留的遗留程序库；
  参见 \zvref{sec:xparse}。
\end{docTcbKey}

\begin{docTcbKey}[library]{external}{}{\mylib{external}}
  为独立文档片段提供外部化支持，
  参见 \zvref{sec:external}。
\end{docTcbKey}

\begin{docTcbKey}[library]{documentation}{}{\mylib{documentation}}
  提供用于排版 \LaTeX\ 文档的额外命令，
  详见 \zvref{sec:documentation}。
\end{docTcbKey}

\begin{docTcbKey}[library]{many}{}{样式，无取值}
  加载 \mylib{skins}、\mylib{breakable}、\mylib{raster}、\mylib{hooks}、
  \mylib{theorems} 和 \mylib{fitting} 程序库。
  如果你想使用 |tcolorbox| 的全部功能，
  但不需要排版代码列表和专门的 \mylib{documentation} 程序库，
  可使用此快捷方式。
\end{docTcbKey}

\begin{docTcbKey}[library]{most}{}{样式，无取值}
  加载除 \mylib{minted}、\mylib{documentation} 和 \mylib{xparse} 之外的所有程序库。
  如果你想使用 |tcolorbox| 的全部功能，
  但不需要使用 \refPkg{minted} 宏包和专门的 \mylib{documentation} 程序库，
  可使用此快捷方式。
\end{docTcbKey}

\begin{docTcbKey}[library]{all}{}{样式，无取值}
  加载所有程序库。仅当你打算使用 \mylib{documentation} 程序库时，
  才使用此快捷方式。
\end{docTcbKey}


\begin{extcolorbox}[runs=2]{intro_packageoverview}
  [title={宏包 \texttt{tcolorbox}},center title,fonttitle=\bfseries,arc=0pt,
  colback=red!10!white,
  interior style={fill tile image*={width=2cm}{pink_marble.png},fill image opacity=0.5},
  colframe=red!50!black]
  \begin{tcolorbox}[beamer,adjusted title=基本功能,colframe=blue!50!black,colback=blue!10!white]
  核心宏包
  \end{tcolorbox}
  \begin{tcbitemize}[raster columns=3,raster before skip=2mm,raster after skip=0pt,
    raster equal height,beamer,colframe=blue!50!black,colback=blue!10!white]
  \tcbitem[adjusted title=高级功能]
    \mylib{breakable}\\
    \mylib{external}\\
    \mylib{fitting}\\
    \mylib{hooks}\\
    \mylib{magazine}\\
    \mylib{poster}\\
    \mylib{raster}\\ 
    \mylib{skins}\\
    \mylib{theorems}\\
    \mylib{vignette}
  \tcbitem[adjusted title=高级代码列表]
    \mylib{listings}\\
    \mylib{listingsutf8}
    \tcblower
    \mylib{minted}
  \tcbitem[adjusted title=文档排版]
    \mylib{documentation}
  \end{tcbitemize}
  \begin{tcolorbox}[beamer,adjusted title=遗留兼容,colframe=blue!50!black,colback=blue!10!white]
    \mylib{xparse}
  \end{tcolorbox}
\end{extcolorbox}
